\mychapter{Conclusões}{cap:conclusion}\lhead{CONCLUSÕES}

Esta dissertação propõe o AIPE como metodologia para selecionar políticas de reposição conforme o contexto da demanda. Sua formulação explicita entradas, cinco etapas, critérios e saídas, distinguindo o procedimento geral das escolhas de cada aplicação. O PSE é o componente de recomendação; associação por perfil e estimação preditiva são suas duas abordagens. Algoritmos de estoque, modelos supervisionados e mecanismos de recuo são instrumentos para instanciá-las.

\section{Respostas às questões de pesquisa}
A primeira questão trata de como estruturar o procedimento. O AIPE encadeia especificação da decisão, representação do contexto, construção de evidência comparável, recomendação e avaliação. Ao separar família e parâmetros, referência e intervenção, desenvolvimento e confirmação, a metodologia explicita quais dados sustentam uma escolha e o que precisa ser observado para julgá-la. O capítulo metodológico define esses compromissos; o desenho experimental mostra sua execução. A existência do procedimento não demonstra, por si, economia operacional.

A segunda questão trata de sua utilidade. Na Bahia, a instanciação por perfil acrescenta 3,89 pontos percentuais de atendimento com aumento pontual de custo de 1,93\%, reduzindo violações de 119 para 100 entre 283 casos com demanda. Na segunda origem, não diferencia a escolha global. No piloto M5, o preditivo original custa 13,9\% mais e atende 4,11 pontos menos que a referência. Sua capacidade de reproduzir rótulos no treinamento não se transfere às lojas excluídas.

A extensão por consequências com recuo apresenta custo médio 1,56\% menor e atendimento 0,27 ponto maior que a global. Entretanto, conserva a referência em 87 de 90 decisões, concentra aproximadamente 92\% da economia em uma série e mantém os mesmos 47 casos com serviço de pelo menos 70\%. Os intervalos exploratórios alcançam diferença zero. O resultado favorece investigar uma intervenção seletiva, sem demonstrar superioridade generalizável.

A terceira questão trata dos limites. Características, rótulos e famílias vencedoras são sensíveis à janela e ao suporte disponível. O piloto contém poucos exemplos por classe e perfis concentrados. Na implementação original, o perfil explícito não altera as recomendações; na floresta por consequências, sua retirada piora o resultado. São evidências condicionais a modelos distintos, não uma conclusão universal sobre perfis.

\section{Contribuição e alcance}
A contribuição principal é o procedimento AIPE e sua instanciação verificável. Ele não pressupõe que a escolha contextual seja sempre preferível: permite examinar quando intervir e quando conservar a referência. O mecanismo de recuo, isoladamente, possui antecedentes; o uso de M5 ou GA--DQN também não constitui novidade própria. Seu alcance decorre da formulação para o problema, da justificativa das decisões e da evidência disponível para cada implementação.

O diagnóstico distingue perdas com causas diferentes. No preditivo original do M5, 21 falhas têm alternativa viável no portfólio, 15 ocorrem sem candidata viável apesar de possibilidade física, 11 decorrem das condições iniciais e quatro casos não têm demanda. Assim, uma melhoria do seletor não resolve toda ruptura. O planejamento perfeito fornece limite condicional, sem representar economia disponível a uma política sem informação futura.

A implementação e os registros computacionais sustentam a rastreabilidade; os resultados brasileiros e públicos avaliam o procedimento em contextos delimitados. Não demonstram implantação contínua, economia financeira realizável ou redução do custo computacional de uso. Origens e sementes repetidas não ampliam a população independente. Custos de cenário, vendas possivelmente censuradas, calibrações diferentes e janelas curtas restringem as conclusões.

\section{Continuidade}
O próximo passo é harmonizar objetivos e calibração e fixar uma avaliação independente da regra de intervenção. O piloto já observado permanece desenvolvimento. Referências globais fortes devem acompanhar ambas as abordagens do PSE, com análise de deteriorações e concentração de benefícios. O fechamento da avaliação deve informar em quais condições recomendar a intervenção, quais deteriorações ela produz e quando conservar a referência.
